\documentclass[../../main.tex]{subfiles}

\begin{document}

\section{Riemann 问题}\label{section:ch11-Riemann问题}

本节详细研究守恒律组
\begin{align}
\mathbf{u}_t+\mathbf{F}(\mathbf{u})_x=0 \quad \text{在 } \mathbb{R}\times(0,\infty) \text{ 中},
\label{align:ch11-守恒律组}
\end{align}
具有分段常数初值
\begin{align}
\mathbf{g}=
\begin{cases}
\mathbf{u}_l & \text{若 } x<0,\\
\mathbf{u}_r & \text{若 } x>0.
\end{cases}
\label{align:ch11-Riemann初值}
\end{align}
这就是\textbf{Riemann 问题}。我们称给定的向量 $\mathbf{u}_l$ 和 $\mathbf{u}_r$ 为左、右初态。

\subsection{简单波}\label{subsection:ch11-简单波}

我们以与 \S11.1.2 大致相同的精神开始研究\eqref{align:ch11-守恒律组}、\eqref{align:ch11-Riemann初值}，即寻找具有特殊形式的解。先前我们搜索行波，即 $\mathbf{u}(x,t)=\mathbf{v}(x-\sigma t)$ 型的解。现在我们寻求\textbf{简单波}。这些是具有结构
\begin{align}
\mathbf{u}(x,t)=\mathbf{v}(w(x,t)) \quad (x\in\mathbb{R},\;t>0),
\label{align:ch11-简单波形式}
\end{align}
的\eqref{align:ch11-守恒律组} 的解，其中 $\mathbf{v}:\mathbb{R}\to\mathbb{R}^m$，$\mathbf{v}=(v_1,\cdots,v_m)$，$w:\mathbb{R}\times[0,\infty)\to\mathbb{R}$ 待求。

为发现 $\mathbf{v}$ 和 $w$ 所需的性质，将\eqref{align:ch11-简单波形式} 代入\eqref{align:ch11-守恒律组}，得到等式
\begin{align}
\dot{\mathbf{v}}(w)w_t+D\mathbf{F}(\mathbf{v}(w))\dot{\mathbf{v}}(w)w_x=0.
\label{align:ch11-简单波代入}
\end{align}
由 \S11.1.2 中方程，$\mathbf{B}=D\mathbf{F}$，我们看到若对某个 $k\in\{1,\cdots,m\}$，$w$ 求解 PDE
\begin{align}
w_t+\lambda_k(\mathbf{v}(w))w_x=0
\label{align:ch11-w方程}
\end{align}
且 $\mathbf{v}$ 求解 ODE
\begin{align}
\dot{\mathbf{v}}(s)=\mathbf{r}_k(\mathbf{v}(s)) \quad (\cdot=\tfrac{\mathrm{d}}{\mathrm{d}s}),
\label{align:ch11-v方程}
\end{align}
则\eqref{align:ch11-简单波代入} 成立。若\eqref{align:ch11-w方程} 和\eqref{align:ch11-v方程} 成立，我们称由\eqref{align:ch11-简单波形式} 定义的函数 $\mathbf{u}$ 为 $k$-简单波。要点在于我们可将\eqref{align:ch11-v方程} 视为向量函数 $\mathbf{v}$ 的 ODE，然后一旦通过求解\eqref{align:ch11-v方程} 找到 $\mathbf{v}$，就可将方程\eqref{align:ch11-w方程} 解释为 $w$ 的标量守恒律。

\begin{definition}\label{definition:ch11-稀疏波曲线}
给定固定状态 $z_0\in\mathbb{R}^m$，我们定义第 $k$ 条\textbf{稀疏波曲线}
\begin{align*}
R_k(z_0)
\end{align*}
为 ODE \eqref{align:ch11-v方程} 的解经过 $z_0$ 在 $\mathbb{R}^m$ 中的路径。
\end{definition}

给定\eqref{align:ch11-v方程} 的解 $\mathbf{v}$，转向 PDE \eqref{align:ch11-w方程}，将其改写为标量守恒律
\begin{align}
w_t+F_k(w)_x=0
\label{align:ch11-标量守恒律}
\end{align}
其中
\begin{align}
F_k(s)\coloneqq\int_0^s\lambda_k(\mathbf{v}(t))\,\mathrm{d}t \quad (s\in\mathbb{R}).
\label{align:ch11-Fk定义}
\end{align}
PDE \eqref{align:ch11-标量守恒律} 将属于 \S3.4 中发展的一般理论，只要 $F_k$ 严格凸（或严格凹）。因此计算
\begin{align}
F_k'(s)=\lambda_k(\mathbf{v}(s)),
\label{align:ch11-Fk一阶导}
\end{align}
\begin{align}
F_k''(s)=D\lambda_k(\mathbf{v}(s))\cdot\dot{\mathbf{v}}(s)=D\lambda_k(\mathbf{v}(s))\cdot\mathbf{r}_k(\mathbf{v}(s)).
\label{align:ch11-Fk二阶导}
\end{align}

由\eqref{align:ch11-Fk二阶导}，函数 $F_k$ 在 $D\lambda_k(z)\cdot\mathbf{r}_k(z)>0$ 时为凸，在 $D\lambda_k(z)\cdot\mathbf{r}_k(z)<0$ 时为凹。函数 $F_k$ 在线性退化情形 $D\lambda_k(z)\cdot\mathbf{r}_k(z)\equiv 0$ 时为线性。这些可能性促使以下定义：

\begin{definition}[i]\label{definition:ch11-真非线性与线性退化} 称对 $(\lambda_k(z),\mathbf{r}_k(z))$ 是\textbf{真非线性的}，若
\begin{align}
D\lambda_k(z)\cdot\mathbf{r}_k(z)\neq 0 \quad \text{对所有 } z\in\mathbb{R}^m.
\label{align:ch11-真非线性}
\end{align}
(ii) 称对 $(\lambda_k(z),\mathbf{r}_k(z))$ 是\textbf{线性退化的}，若
\begin{align}
D\lambda_k(z)\cdot\mathbf{r}_k(z)=0 \quad \text{对所有 } z\in\mathbb{R}^m.
\label{align:ch11-线性退化}
\end{align}
\end{definition}

\textbf{记号。} 若对 $(\lambda_k,\mathbf{r}_k)$ 是真非线性的，写
\begin{align*}
R_k^+(z_0)\coloneqq\{z\in R_k(z_0)\mid\lambda_k(z)>\lambda_k(z_0)\}
\end{align*}
和
\begin{align*}
R_k^-(z_0)\coloneqq\{z\in R_k(z_0)\mid\lambda_k(z)<\lambda_k(z_0)\}.
\end{align*}
则 $R_k(z_0)=R_k^+(z_0)\cup\{z_0\}\cup R_k^-(z_0)$。

\subsection{稀疏波}\label{subsection:ch11-稀疏波}

我们再次转向 Riemann 问题\eqref{align:ch11-守恒律组}、\eqref{align:ch11-Riemann初值}。

\begin{theorem}[$k$-稀疏波的存在性。]\label{theorem:ch11-k稀疏波存在性}假设对某个 $k\in\{1,\cdots,m\}$，
\begin{enumerate}[(i)]
\item 对 $(\lambda_k,\mathbf{r}_k)$ 是真非线性的，且
\item $\mathbf{u}_r\in R_k^+(\mathbf{u}_l)$。
\end{enumerate}
则存在 Riemann 问题\eqref{align:ch11-守恒律组}、\eqref{align:ch11-Riemann初值} 的连续积分解 $\mathbf{u}$，它是沿过原点直线为常数的 $k$-简单波。我们称 $\mathbf{u}$ 为（中心）$k$-稀疏波。
\end{theorem}

\begin{proof}
首先选取 $w_l,w_r\in\mathbb{R}$ 使得 $\mathbf{u}_l=\mathbf{v}(w_l)$，$\mathbf{u}_r=\mathbf{v}(w_r)$。暂时假设 $w_l<w_r$。考虑由 PDE \eqref{align:ch11-标量守恒律} 连同初值条件
\begin{align*}
g=\begin{cases}
w_l & \text{若 } x<0,\\
w_r & \text{若 } x>0
\end{cases}
\end{align*}
组成的标量 Riemann 问题。由假设 (ii)，$\lambda_k(\mathbf{u}_r)>\lambda_k(\mathbf{u}_l)$，即由\eqref{align:ch11-Fk一阶导}，$F_k'(w_r)>F_k'(w_l)$。由 (i)，由\eqref{align:ch11-Fk定义} 定义的函数 $F_k$ 严格凸。因此可对该标量 Riemann 问题应用 \S3.4.4\cref{theorem:ch03-riemann}，其唯一弱解是连接状态 $w_l$ 和 $w_r$ 的连续稀疏波。具体地，
\begin{align*}
w(x,t)=
\begin{cases}
w_l & \text{若 } \dfrac{x}{t}<F_k'(w_l),\\
G_k\left(\dfrac{x}{t}\right) & \text{若 } F_k'(w_l)<\dfrac{x}{t}<F_k'(w_r),\\
w_r & \text{若 } \dfrac{x}{t}>F_k'(w_r),
\end{cases}
\end{align*}
其中 $G_k=(F_k')^{-1}$。因此 $\mathbf{u}(x,t)=\mathbf{v}(w(x,t))$，其中 $\mathbf{v}$ 求解 ODE \eqref{align:ch11-v方程} 并经过 $\mathbf{u}_l$，是\eqref{align:ch11-守恒律组}、\eqref{align:ch11-Riemann初值} 的连续积分解。

$w_l>w_r$ 的情形类似处理，因为此时 $F_k$ 为凹。
\end{proof}

\subsection{激波，接触间断}\label{subsection:ch11-激波接触间断}

接下来考虑状态 $\mathbf{u}_l$ 和 $\mathbf{u}_r$ 可能不由上述稀疏波连接、而由激波连接的可能性。

\subsubsection*{激波集。} 回忆 \S11.1.1 的 Rankine--Hugoniot 条件，我们看到要存在这样的激波，必须有等式 $\mathbf{F}(\mathbf{u}_l)-\mathbf{F}(\mathbf{u}_r)=\sigma(\mathbf{u}_l-\mathbf{u}_r)$，其中 $\sigma\in\mathbb{R}$。此观察激发以下定义：

\begin{definition}\label{definition:ch11-激波集}
给定固定状态 $z_0\in\mathbb{R}^m$，我们定义\textbf{激波集}
\begin{align*}
S(z_0)\coloneqq\{z\in\mathbb{R}^m\mid\mathbf{F}(z)-\mathbf{F}(z_0)=\sigma(z-z_0)\text{ 对某常数 }\sigma=\sigma(z,z_0)\}.
\end{align*}
\end{definition}

\begin{theorem}[激波集的结构。]\label{theorem:ch11-激波集结构}固定 $z_0\in\mathbb{R}^m$。在 $z_0$ 的某个邻域内，$S(z_0)$ 由 $m$ 条光滑曲线 $S_k(z_0)$（$k=1,\cdots,m$）的并组成，具有以下性质：
\begin{enumerate}[(i)]
\item 曲线 $S_k(z_0)$ 经过 $z_0$，切向量为 $\mathbf{r}_k(z_0)$。
\item $\displaystyle\lim_{\substack{z\to z_0\\ z\in S_k(z_0)}}\sigma(z,z_0)=\lambda_k(z_0)$。
\item $\sigma(z,z_0)=\dfrac{\lambda_k(z)+\lambda_k(z_0)}{2}+O(|z-z_0|^2)$，当 $z\to z_0$ 且 $z\in S_k(z_0)$ 时。
\end{enumerate}
\end{theorem}

\begin{proof}
1. 定义
\begin{align*}
\mathbf{B}(z)\coloneqq\int_0^1 D\mathbf{F}(z_0+t(z-z_0))\,\mathrm{d}t \quad (z\in\mathbb{R}^m).
\end{align*}
则 $\mathbf{B}(z)(z-z_0)=\mathbf{F}(z)-\mathbf{F}(z_0)$。特别地，$z\in S(z_0)$ 当且仅当
\begin{align}
(\mathbf{B}(z)-\sigma I)(z-z_0)=0
\label{align:ch11-激波集方程}
\end{align}
对某个标量 $\sigma=\sigma(z,z_0)$。

2. 注意 $\mathbf{B}(z_0)=D\mathbf{F}(z_0)$。由严格双曲性，特征多项式 $\lambda\mapsto\det(\lambda I-\mathbf{B}(z_0))$ 有 $m$ 个互异实根，故当 $z$ 接近 $z_0$ 时，多项式 $\lambda\mapsto\det(\lambda I-\mathbf{B}(z))$ 同样有 $m$ 个互异根。由 \S11.1.2\cref{theorem:ch11-特征值光滑依赖}，在 $z_0$ 附近存在光滑函数 $\widehat{\lambda}_1(z)<\cdots<\widehat{\lambda}_m(z)$ 和单位向量 $\{\widehat{\mathbf{r}}_k(z),\widehat{\mathbf{l}}_k(z)\}_{k=1}^m$ 满足
\begin{align*}
\widehat{\lambda}_k(z_0)=\lambda_k(z_0),\;\widehat{\mathbf{r}}_k(z_0)=\mathbf{r}_k(z_0),\;\widehat{\mathbf{l}}_k(z_0)=\mathbf{l}_k(z_0)
\end{align*}
和
\begin{align*}
\begin{cases}
\mathbf{B}(z)\widehat{\mathbf{r}}_k(z)=\widehat{\lambda}_k(z)\widehat{\mathbf{r}}_k(z),\\
\widehat{\mathbf{l}}_k(z)\mathbf{B}(z)=\widehat{\lambda}_k(z)\widehat{\mathbf{l}}_k(z).
\end{cases}
\end{align*}

3. 方程\eqref{align:ch11-激波集方程} 成立，只要 $\sigma=\widehat{\lambda}_k(z)$ 对某个 $k\in\{1,\cdots,m\}$ 且 $(z-z_0)$ 平行于 $\widehat{\mathbf{r}}_k(z)$。这些条件等价于 $\widehat{\mathbf{l}}_l(z)\cdot(z-z_0)=0$（$l\neq k$）。这些等式构成 $m-1$ 个方程、$m$ 个未知量。定义 $\boldsymbol{\Phi}_k:\mathbb{R}^m\to\mathbb{R}^{m-1}$ 为
\begin{align*}
\boldsymbol{\Phi}_k(z)\coloneqq(\cdots,\widehat{\mathbf{l}}_{k-1}(z)\cdot(z-z_0),\widehat{\mathbf{l}}_{k+1}(z)\cdot(z-z_0),\cdots).
\end{align*}
由于 $\{\mathbf{l}_k(z_0)\}_{k=1}^m$ 构成 $\mathbb{R}^m$ 的基，$\operatorname{rank}D\boldsymbol{\Phi}_k(z_0)=m-1$。由隐函数定理，存在光滑曲线 $\boldsymbol{\phi}_k:\mathbb{R}\to\mathbb{R}^m$ 使得 $\boldsymbol{\phi}_k(0)=z_0$ 且 $\boldsymbol{\Phi}_k(\boldsymbol{\phi}_k(t))=0$。曲线 $\boldsymbol{\phi}_k(\cdot)$ 的路径定义 $S_k(z_0)$。

4. 我们有 $\boldsymbol{\phi}_k(t)=z_0+\mu(t)\widehat{\mathbf{r}}_k(\boldsymbol{\phi}_k(t))$，其中 $\mu(0)=0$，$\dot{\mu}(0)=1$。对 $t$ 求导并令 $t=0$，得 $\dot{\boldsymbol{\phi}}_k(0)=\mathbf{r}_k(z_0)$。断言 (i) 得证。

5. 存在光滑函数 $\sigma:\mathbb{R}^m\times\mathbb{R}^m\to\mathbb{R}$ 使得 $\mathbf{F}(\boldsymbol{\phi}_k(t))-\mathbf{F}(z_0)=\sigma(\boldsymbol{\phi}_k(t),z_0)(\boldsymbol{\phi}_k(t)-z_0)$。对 $t$ 求导并令 $t=0$，得 $D\mathbf{F}(z_0)\dot{\boldsymbol{\phi}}_k(0)=\sigma(z_0,z_0)\dot{\boldsymbol{\phi}}_k(0)$。由 $\dot{\boldsymbol{\phi}}_k(0)=\mathbf{r}_k(z_0)$，$\sigma(z_0,z_0)=\lambda_k(z_0)$。断言 (ii) 得证。

6--7. 对等式两次求导并与稀疏波曲线的参数化比较，取与 $\mathbf{l}_k(z_0)$ 的点积，得 $2\dot{\sigma}(0)=\dot{\lambda}_k(0)$。由此推出 $2\sigma(t)-\lambda_k(z_0)-\lambda_k(t)=O(t^2)$。断言 (iii) 得证。
\end{proof}

由\cref{theorem:ch11-激波集结构}(iii)，曲线 $R_k(z_0)$ 和 $S_k(z_0)$ 在 $z_0$ 处至少一阶一致。以下断言在线性退化情形下这些曲线实际上重合。

\begin{theorem}[线性退化。]\label{theorem:ch11-线性退化}假设对某个 $k\in\{1,\cdots,m\}$，对 $(\lambda_k,\mathbf{r}_k)$ 是线性退化的。则对每个 $z_0\in\mathbb{R}^m$，
\begin{enumerate}[(i)]
\item $R_k(z_0)=S_k(z_0)$，且
\item $\sigma(z,z_0)=\lambda_k(z)=\lambda_k(z_0)$ 对所有 $z\in S_k(z_0)$。
\end{enumerate}
\end{theorem}

\begin{proof}
令 $\mathbf{v}=\mathbf{v}(s)$ 求解 ODE
\begin{align*}
\begin{cases}
\dot{\mathbf{v}}(s)=\mathbf{r}_k(\mathbf{v}(s)) & (s\in\mathbb{R}),\\
\mathbf{v}(0)=z_0.
\end{cases}
\end{align*}
则映射 $s\mapsto\lambda_k(\mathbf{v}(s))$ 为常数，故
\begin{align*}
\mathbf{F}(\mathbf{v}(s))-\mathbf{F}(z_0)=\int_0^s D\mathbf{F}(\mathbf{v}(t))\dot{\mathbf{v}}(t)\,\mathrm{d}t=\int_0^s \lambda_k(\mathbf{v}(t))\mathbf{r}_k(\mathbf{v}(t))\,\mathrm{d}t=\lambda_k(z_0)(\mathbf{v}(s)-z_0).
\end{align*}
\end{proof}

\subsubsection*{接触间断，激波。} 我们接下来根据\cref{theorem:ch11-激波集结构} 和\cref{theorem:ch11-线性退化} 分析通过某种激波连接两个给定状态 $\mathbf{u}_l$ 和 $\mathbf{u}_r$ 来求解 Riemann 问题的可能性。

\textbf{接触间断。} 首先假设 $(\lambda_k,\mathbf{r}_k)$ 是线性退化的且 $\mathbf{u}_r\in S_k(\mathbf{u}_l)$。我们定义积分解
\begin{align}
\mathbf{u}(x,t)=
\begin{cases}
\mathbf{u}_l & \text{若 } x<\sigma t,\\
\mathbf{u}_r & \text{若 } x>\sigma t,
\end{cases}
\label{align:ch11-接触间断解}
\end{align}
其中 $\sigma=\sigma(\mathbf{u}_r,\mathbf{u}_l)=\lambda_k(\mathbf{u}_l)=\lambda_k(\mathbf{u}_r)$。由于 $\lambda_k(\mathbf{u}_l)=\lambda_k(\mathbf{u}_r)=\sigma$，左侧和右侧的投影特征线平行于间断线。我们物理解释为流体粒子不穿过间断。直线 $x=\sigma t$ 称为 $k$-接触间断。

\textbf{激波。} 接下来转向 $(\lambda_k,\mathbf{r}_k)$ 真非线性且 $\mathbf{u}_r\in S_k(\mathbf{u}_l)$ 的情形。考虑积分解
\begin{align}
\mathbf{u}(x,t)=
\begin{cases}
\mathbf{u}_l & \text{若 } x<\sigma t,\\
\mathbf{u}_r & \text{若 } x>\sigma t,
\end{cases}
\label{align:ch11-激波解}
\end{align}
其中 $\sigma=\sigma(\mathbf{u}_r,\mathbf{u}_l)$。我们看到有两种本质不同的情形：$\lambda_k(\mathbf{u}_r)<\lambda_k(\mathbf{u}_l)$ 或 $\lambda_k(\mathbf{u}_l)<\lambda_k(\mathbf{u}_r)$。由\cref{theorem:ch11-激波集结构}(iii)，当 $\mathbf{u}_r$ 充分接近 $\mathbf{u}_l$ 时，有
\begin{align}
\lambda_k(\mathbf{u}_r)<\sigma(\mathbf{u}_r,\mathbf{u}_l)<\lambda_k(\mathbf{u}_l)
\label{align:ch11-Lax熵条件}
\end{align}
或反向不等式。类比 \S3.4 中的熵条件，我们同意拒绝后一种不等式所允许的``非物理激波''。我们取\eqref{align:ch11-Lax熵条件} 为物理正确的。

\begin{definition}\label{definition:ch11-Lax熵条件}
假设对 $(\lambda_k,\mathbf{r}_k)$ 在 $\mathbf{u}_l$ 处真非线性。称对 $(\mathbf{u}_l,\mathbf{u}_r)$ 是\textbf{容许的}，若
\begin{align*}
\mathbf{u}_r\in S_k(\mathbf{u}_l)
\end{align*}
且
\begin{align}
\lambda_k(\mathbf{u}_r)<\sigma(\mathbf{u}_r,\mathbf{u}_l)<\lambda_k(\mathbf{u}_l).
\label{align:ch11-Lax条件}
\end{align}
我们称\eqref{align:ch11-Lax条件} 为\textbf{Lax 熵条件}。若 $(\mathbf{u}_l,\mathbf{u}_r)$ 是容许的，称由\eqref{align:ch11-激波解} 定义的解 $\mathbf{u}$ 为 $k$-激波。
\end{definition}

\textbf{记号。} 若对 $(\lambda_k,\mathbf{r}_k)$ 真非线性，写
\begin{align*}
S_k^+(z_0)\coloneqq\{z\in S_k(z_0)\mid\lambda_k(z_0)<\sigma(z,z_0)<\lambda_k(z)\}
\end{align*}
和
\begin{align*}
S_k^-(z_0)\coloneqq\{z\in S_k(z_0)\mid\lambda_k(z)<\sigma(z,z_0)<\lambda_k(z_0)\}.
\end{align*}
则在 $z_0$ 附近 $S_k(z_0)=S_k^+(z_0)\cup\{z_0\}\cup S_k^-(z_0)$。注意对 $(\mathbf{u}_l,\mathbf{u}_r)$ 容许当且仅当 $\mathbf{u}_r\in S_k^-(\mathbf{u}_l)$。

\subsection{Riemann 问题的局部解}\label{subsection:ch11-Riemann局部解}

接下来我们将稀疏波和激波曲线的物理相关部分粘合在一起。

\begin{definition}[i]\label{definition:ch11-T曲线} 若对 $(\lambda_k,\mathbf{r}_k)$ 真非线性，写
\begin{align*}
T_k(z_0)\coloneqq R_k^+(z_0)\cup\{z_0\}\cup S_k^-(z_0).
\end{align*}
(ii) 若对 $(\lambda_k,\mathbf{r}_k)$ 线性退化，设
\begin{align*}
T_k(z_0)\coloneqq R_k(z_0)=S_k(z_0).
\end{align*}
\end{definition}

由\cref{theorem:ch11-激波集结构}(ii)，曲线 $T_k(z_0)$ 是 $C^1$ 的。我们看到附近状态 $\mathbf{u}_l$ 和 $\mathbf{u}_r$ 可由 $k$-稀疏波、激波或接触间断连接，只要 $\mathbf{u}_l\in T_k(\mathbf{u}_r)$。

\begin{theorem}[Riemann 问题的局部解。]\label{theorem:ch11-Riemann局部解}假设对每个 $k=1,\cdots,m$，对 $(\lambda_k,\mathbf{r}_k)$ 是真非线性的或线性退化的。进一步假设左状态 $\mathbf{u}_l$ 给定。则对每个充分接近 $\mathbf{u}_l$ 的右状态 $\mathbf{u}_r$，存在 Riemann 问题的积分解 $\mathbf{u}$，它沿过原点的直线为常数。
\end{theorem}

\begin{proof}
1. 我们将逆函数定理应用于映射 $\Phi:\mathbb{R}^m\to\mathbb{R}^m$，在 0 附近定义如下。对每个曲线族 $T_k$（$k=1,\cdots,m$），选择非奇异参数 $\tau_k$ 测量弧长。

2. 给定 $t=(t_1,\cdots,t_m)\in\mathbb{R}^m$，$|t|$ 小，定义 $\Phi(t)=z$。先写 $\mathbf{u}_l=z_0$。然后选择状态 $z_1,\cdots,z_m$ 满足
\begin{align*}
\begin{cases}
z_1\in T_1(z_0),\;\tau_1(z_1)-\tau_1(z_0)=t_1,\\
z_2\in T_2(z_1),\;\tau_2(z_2)-\tau_2(z_1)=t_2,\\
\quad\vdots\\
z_m\in T_m(z_{m-1}),\;\tau_m(z_m)-\tau_m(z_{m-1})=t_m.
\end{cases}
\end{align*}
写 $z=z_m$，定义 $\Phi(t)=z$。注意 $\Phi$ 是 $C^1$ 的且 $\Phi(0)=z_0$。

3. 我们断言 $D\Phi(0)$ 非奇异。事实上，
\begin{align*}
\Phi(0,\cdots,t_k,\cdots,0)-\Phi(0,\cdots,0)=t_k\mathbf{r}_k(z_0)+o(t_k)
\end{align*}
故 $\dfrac{\partial\Phi}{\partial t_k}(0)=\mathbf{r}_k(z_0)$，从而
\begin{align*}
D\Phi(0)=(\mathbf{r}_1(z_0),\cdots,\mathbf{r}_m(z_0))_{m\times m}.
\end{align*}
此矩阵非奇异，因为 $\{\mathbf{r}_k(z_0)\}_{k=1}^m$ 是基。

4. 由逆函数定理，对每个充分接近 $\mathbf{u}_l$ 的状态 $\mathbf{u}_r$，存在唯一参数 $t=(t_1,\cdots,t_m)$ 接近零使得 $\Phi(t)=\mathbf{u}_r$。

若 $z_{k-1}$ 和 $z_k$ 由 $k$-稀疏波连接，该波为
\begin{align*}
\begin{cases}
z_{k-1} & \text{若 } \dfrac{x}{t}<\lambda_k(z_{k-1}),\\
G_k\left(\dfrac{x}{t}\right) & \text{若 } \lambda_k(z_{k-1})<\dfrac{x}{t}<\lambda_k(z_k),\\
z_k & \text{若 } \lambda_k(z_k)<\dfrac{x}{t},
\end{cases}
\end{align*}
其中 $G_k=(F_k')^{-1}$。若 $z_{k-1},z_k$ 由 $k$-激波连接，其形式为
\begin{align*}
\begin{cases}
z_{k-1} & \text{若 } \dfrac{x}{t}<\sigma(z_k,z_{k-1}),\\
z_k & \text{若 } \sigma(z_k,z_{k-1})<\dfrac{x}{t}.
\end{cases}
\end{align*}
其中 $\lambda_k(z_k)<\sigma(z_k,z_{k-1})<\lambda_k(z_{k-1})$。在两种情形中波在区域 $\lambda_k(z_0)-\varepsilon<\dfrac{x}{t}<\lambda_k(z_0)+\varepsilon$ 外为常数，对小 $\varepsilon>0$，只要 $z_k,z_{k-1}$ 充分接近 $z_0$。这对 $k=1,\cdots,m$ 成立。

由于 $\lambda_1(z_0)<\cdots<\lambda_m(z_0)$，我们看到连接 $\mathbf{u}_l=z_0$ 到 $z_1$、$z_1$ 到 $z_2$、$z_2$ 到 $z_3$、$\cdots$、$z_{m-1}$ 到 $z_m=\mathbf{u}_r$ 的稀疏波、激波和/或接触间断不相交。
\end{proof}

\end{document}